% Section 1.9: Board Support Packages (BSP)

\section{Creating Board Support Packages (BSP)}
\label{sec:bsp-implementation}

A Board Support Package bridges the gap between the Hardware Abstraction Layer (which is chip-specific) and your application (which should be board-agnostic). The BSP provides board-specific configurations, pin mappings, and convenience functions for the exact hardware you're using.

\subsection{BSP vs HAL: The Distinction}
\label{subsec:bsp-vs-hal}

\begin{table}[H]
    \centering
    \begin{tabular}{l l}
        \toprule
        \textbf{HAL (Hardware Abstraction)} & \textbf{BSP (Board Support)} \\\\\midrule
        Chip-specific (EFR32MG24 series) & Board-specific (XIAO MG24 Sense) \\\\
        133 peripherals available & Only exposed board features \\\\
        Generic pin API & Specific pin functions (LED, button, etc.) \\\\
        No board assumptions & Assumes specific PCB layout \\\\
        Reusable across boards & Board revision specific \\\\
bottomrule
    \end{tabular}
    \caption{HAL vs BSP responsibilities}
    \label{tab:hal-vs-bsp}
\end{table}

\subsection{XIAO MG24 Sense Board Layout}
\label{subsec:xiao-layout}

The Seeed Studio XIAO MG24 Sense includes:

\begin{itemize}
    \item EFR32MG24B220F1536IM48 microcontroller
    \item Two user LEDs (red and blue)
    \item One user button
    \item I2C connector (QWIIC/Stemma QT compatible)
    \item 2.4 GHz antenna
    \item USB-C connector (power + debug)
    \item SWD debug pins exposed
\end{itemize}

\begin{figure}[H]
    \centering
    \includegraphics[width=0.8\textwidth]{xiao-mg24-pinout}
    \caption{XIAO MG24 Sense pinout diagram}
    \label{fig:xiao-pinout}
\end{figure}

\subsection{BSP Structure}
\label{subsec:bsp-structure}

A typical BSP crates:

\begin{minted}{text}
xiao-mg24-bsp/
├── Cargo.toml
├── README.md
└── src/
    ├── lib.rs              # Board initialization
    ├── led.rs              # LED abstraction
    ├── button.rs           # Button abstraction
    ├── sensors.rs          # Onboard sensors
    ├── debug.rs            # Debug utilities
    └── pins.rs             # Pin mappings
\end{minted}

\subsection{Implementing LED Abstraction}
\label{subsec:led-abstraction}

\begin{minted}{rust}
//! src/led.rs
use efr32mg24_hal::gpio::{Output, Pin, PushPull};

pub struct UserLeds {
    pub red: Pin<PA, PIN5, Output<PushPull>>,
    pub blue: Pin<PA, PIN6, Output<PushPull>>,
}

impl UserLeds {
    pub fn new(gpio: GPIO) -> Self {
        let red = gpio.pa5.into_output();
        let blue = gpio.pa6.into_output();
        
        Self { red, blue }
    }
    
    pub fn both_off(&mut self) {
        self.red.set_low().unwrap();
        self.blue.set_low().unwrap();
    }
}
\end{minted}

\subsection{Pin Mappings as Type Aliases}
\label{subsec:pin-mappings}

\begin{minted}{rust}
//! src/pins.rs
use efr32mg24_hal::gpio::{Input, Pin, PullUp};

type ButtonPin = Pin<PA, PIN7, Input<PullUp>>;

pub struct Buttons {
    pub user: ButtonPin,
}

impl Buttons {
    pub fn new(gpio: GPIO) -> Self {
        let user = gpio.pa7.into_input().pull_up();
        Self { user }
    }
    
    pub fn is_pressed(&self) -> bool {
        self.user.is_low().unwrap()
    }
}
\end{minted}

\subsection{Board Initialization}
\label{subsec:board-init}

\begin{minted}{rust}
//! src/lib.rs
pub struct Board {
    pub leds: UserLeds,
    pub buttons: Buttons,
    pub i2c: I2c0,
}

impl Board {
    pub fn init() -> Self {
        // Critical: Enable GPIO clock
        cmu.enable_gpio_clock();
        
        let gpio = GPIO::new();
        let leds = UserLeds::new(gpio);
        let buttons = Buttons::new(gpio);
        let i2c = i2c::I2c0::init();
        
        Board { leds, buttons, i2c }
    }
}
\end{minted}

\begin{example}
Usage in application code:
\end{example}

\begin{minted}{rust}
use xiao_mg24_bsp::Board;

#[embassy_executor::task]
async fn blinky() {
    let mut board = Board::init();
    
    loop {
        board.leds.red.toggle().unwrap();
        Timer::after(Duration::from_millis(500)).await;
    }
}
\end{minted}

\begin{keytakeaways}
    \item BSP provides board-specific abstractions over HAL
    \item BSP is board revision specific, HAL is chip family specific
    \item Pin mappings should be types, not raw numbers
    \item Board initialization is critical and should be explicit
    \item BSP enables portable application code
\end{keytakeaways}

\begin{exercises}
    \item\textbf{Easy}: Create a BSP for a different board (e.g., SparkFun Thing Plus).
    
    \item\textbf{Medium}: Add driver support for an I2C sensor on the XIAO MG24.
    
    \item\textbf{Hard}: Design a BSP that supports multiple board revisions through feature flags.
\end{exercises}
